% =============================================================================
% 2-claims template — standalone-compilable lifecycle stage contract
%
% Structure: Editor's Chair → RQs → Matrix → Detail → Alignment
% Format: 1 PRIMARY + N SUPPORTING, venue-coupled
%
% The editor's chair question works TWICE:
%   top-down: generates RQs (if an RQ doesn't help answer the editor, drop it)
%   bottom-up: validates claims (if a claim can't produce a one-sentence answer, wrong venue)
% =============================================================================
\documentclass[11pt]{article}
\usepackage[margin=1in]{geometry}
\usepackage{parskip}
\usepackage{booktabs}
\usepackage{tabularx}
\usepackage{hyperref}
\usepackage{xcolor}
\newcommand{\status}[1]{\textsc{#1}}
\newcommand{\supported}{\textcolor{green!50!black}{\status{supported}}}
\newcommand{\weak}{\textcolor{orange!80!black}{\status{weak}}}
\newcommand{\gap}{\textcolor{red!70!black}{\status{gap}}}
\newcommand{\pending}{\textcolor{blue!60!black}{\status{pending}}}
\newcommand{\needprobe}[1]{\textcolor{red}{\textbf{[NEED PROBE]} #1}}
\title{Claims: <Paper Title> (<Venue>)}
\date{YYYY-MM-DD}
\begin{document}
\maketitle

% ─────────────────────────────────────────────────────────────────────────────
% 1. EDITOR'S CHAIR TEST
%    Source: _venue/playbook-<venue>
%    Every primary claim must have a one-sentence answer.
% ─────────────────────────────────────────────────────────────────────────────
\section*{Editor's Chair Test}

\textbf{<Venue>:} ``<editor's chair question from _venue/playbook>''

Every primary claim must have a one-sentence answer to this question.

% ─────────────────────────────────────────────────────────────────────────────
% 2. RESEARCH QUESTIONS
%    Venue-shaped: the wording depends on what the target editor rewards.
%    RQs live here (claims), not in the pitch (pitch is venue-independent).
% ─────────────────────────────────────────────────────────────────────────────
\section*{Research Questions}

\begin{enumerate}
\item[\textbf{RQ1.}] <question, shaped by what the venue rewards>
\item[\textbf{RQ2.}] <question>
\item[\textbf{RQ3.}] <question>
\end{enumerate}

\bigskip

% RQ-to-Claim mapping: the skeleton of the Results section.
% "Why this RQ" column makes the venue→RQ generation link explicit.
\begin{small}
\begin{tabularx}{\textwidth}{@{}p{0.05\textwidth}p{0.05\textwidth}X p{0.28\textwidth}@{}}
\toprule
RQ & Claims & Mapping & Why this RQ for this venue \\
\midrule
RQ1 & C1, C2 & <how the claims answer this RQ> & <why the editor's chair question led to this RQ> \\
RQ2 & C3     & <how the claims answer this RQ> & <why> \\
RQ3 & C4     & <how the claims answer this RQ> & <why> \\
\bottomrule
\end{tabularx}
\end{small}

% ─────────────────────────────────────────────────────────────────────────────
% 3. CLAIM-EVIDENCE MATRIX
%    Index: one row per claim, status at a glance.
%    One claim marked [primary]; rest are support.
% ─────────────────────────────────────────────────────────────────────────────
\section*{Claim-Evidence Matrix}

% =========================================================
% Para [claims.ledger] Result -- claim / status at a glance
% =========================================================
\begin{small}
\begin{tabularx}{\textwidth}{@{}p{0.06\textwidth}p{0.08\textwidth}X p{0.12\textwidth}@{}}
\toprule
ID & Role & Claim & Status \\
\midrule
C1 & \textbf{[primary]} & <claim, as the paper states it> & \supported \\
\addlinespace
C2 & support & <claim> & \weak \\
\addlinespace
C3 & support & <claim> & \gap \\
\bottomrule
\end{tabularx}
\end{small}

% ─────────────────────────────────────────────────────────────────────────────
% 4. PER-CLAIM DETAIL
%    One \subsection* per claim. Four slots per paragraph:
%    S1 = claim + verdict, S2 = verified statistic, S3 = interpretation, S4 = source + caveat
% ─────────────────────────────────────────────────────────────────────────────
\section*{Per-Claim Detail}

\subsection*{C1. <short title> [primary] (\supported)}

% =========================================================
% Para [claims.c1] Result -- <the point>
% =========================================================
%% ---- P1.S1 ----
<claim statement>.
%
%% ---- P1.S2 ----
<verified statistic, spec, N>.
%
%% ---- P1.S3 ----
<one-line interpretation>.
%
%% ---- P1.S4 ----
Source: <real file>; <caveat if any>.

% ... more claims: C2, C3, etc. ...
% For weak/GAP claims: state the gap and the route instead of a statistic.

% ─────────────────────────────────────────────────────────────────────────────
% 5. DISCUSSION-ONLY INTERPRETATIONS (optional)
%    Interpretive findings that belong in Discussion, not Results.
%    Explicitly marked to prevent creep into Results claims.
% ─────────────────────────────────────────────────────────────────────────────
\section*{Discussion-Only Interpretations}

\textbf{D1. <label> (interpretive):} <description>. No direct measurement; interpretive frame only.

% ─────────────────────────────────────────────────────────────────────────────
% 6. ROBUSTNESS (optional)
%    Reported as a design strength in Methods, not claimed as a finding.
% ─────────────────────────────────────────────────────────────────────────────
\section*{Robustness (reported in Methods, not claimed)}

<sensitivity analyses: clustering, alt specs, multiple testing, exclusion criteria>

% ─────────────────────────────────────────────────────────────────────────────
% 7. PENDING EVIDENCE
%    Probes/tasks not yet run. Dependency chain. Backup venue if key evidence weak.
% ─────────────────────────────────────────────────────────────────────────────
\section*{Pending Evidence}

\textbf{<probe id> (<label>, <status>):} <what it will test>. <which claim it upgrades>.

\textbf{Backup}: If <key probe> is weak $\to$ <fallback venue> with <adjusted claim>.

\textbf{Exploratory (supplement):} <analyses held from main claims>.

% ─────────────────────────────────────────────────────────────────────────────
% 8. EDITOR'S CHAIR ALIGNMENT (at the bottom — validates the whole ledger)
%    Three-column table: RQ → Claims → Editor's Chair Answer
%    Merged with Venue Fit: scale, design strength, risks.
%    Diagnostic: any claim missing a column has a structural problem.
% ─────────────────────────────────────────────────────────────────────────────
\section*{Editor's Chair Alignment}

\textbf{Editor's chair question (<Venue>):} ``<question>''

The editor's chair question works twice: once at the top (generating RQs ---
if an RQ does not help answer this question, drop it) and once at the bottom
(validating claims --- if a claim cannot produce a one-sentence answer, it is
either the wrong claim or the wrong venue).

\bigskip

\begin{small}
\begin{tabularx}{\textwidth}{@{}p{0.05\textwidth} p{0.26\textwidth} p{0.13\textwidth} X@{}}
\toprule
RQ & Research Question & Claims & Editor's Chair Answer \\
\midrule
RQ1 & <question> & C1, C2 & ``<one-sentence answer to the editor>'' \\
\addlinespace
RQ2 & <question> & C3 & ``<answer>'' \\
\addlinespace
RQ3 & <question> & C4 & ``<answer>'' \\
\bottomrule
\end{tabularx}
\end{small}

\bigskip

\textbf{Venue fit:}
\begin{itemize}
\item \textbf{Scale}: <dataset size, design, what makes this competitive>
\item \textbf{Strength}: <what the venue rewards that this paper delivers>
\item \textbf{Risk}: <where reviewers may push back>
\end{itemize}

\textbf{Diagnostic:} Any claim that cannot fill all three columns has a problem.
\begin{itemize}
\item No RQ $\to$ orphan finding, not driven by a question.
\item No editor's chair answer $\to$ wrong claim or wrong venue.
\item RQ without a claim $\to$ unanswered question (evidence is GAP or aspirational).
\end{itemize}

\end{document}
